\documentclass[10pt,a4paper]{amsart}
\usepackage[margin=19mm]{geometry}
\usepackage{amsmath,amssymb,mathtools}
\usepackage[scr=rsfs,cal=euler]{mathalfa}
\usepackage{xfrac}
\usepackage[unicode,hidelinks,bookmarks=false]{hyperref}
\newcommand{\ie}{i.e.\ }
\newcommand{\cf}{cf.\ }
\newcommand{\resp}{respectively\ }
\newcommand{\Subsection}{\subsection}
\providecommand{\nobreakdash}{\nobreak\hbox{-}}
\setlength{\emergencystretch}{2em}
\multlinegap0pt
\usepackage{txfonts}
\usepackage{eucal}
\usepackage{array,longtable}
\usepackage{float}
\usepackage{url}
\usepackage{appendix}
\usepackage{enumitem}
\usepackage{algorithm}
\usepackage{algpseudocode}
\newtheorem{theorem}{Theorem}[section]
\newtheorem{corollary}[theorem]{Corollary}
\newtheorem{lemma}[theorem]{Lemma}
\newtheorem{proposition}[theorem]{Proposition}
\newtheorem{definition}[theorem]{Definition}
\newtheorem{notation}[theorem]{Notation}
\newtheorem{example}[theorem]{Example}
\newtheorem{problem}[theorem]{Problem}
\newtheorem{question}[theorem]{Question}
\newtheorem{conjecture}{Conjecture}
\newtheorem{claim}{Claim}
\numberwithin{equation}{section}
\newtheorem{remark}[theorem]{Remark}
\newtheorem{conclusion}[theorem]{Conclusion}
\newcounter{condition}
\newcommand{\conditionheading}[1]{%
  \refstepcounter{condition}%
  \subsubsection*{\textbf{\thecondition}}%
  \label{#1}%
}
\allowdisplaybreaks[4]
\begin{document}
\title[Substitution and quotient of the isotropy group action]{Substitution and quotient of the isotropy group action}
\author{Xin Li, Yu Wang; Shenglong Hu}
\date{}
\maketitle
\noindent\emph{Source and licence.} Substitution and quotient of the isotropy group action. Version arXiv:2607.15069v2 (CC BY 4.0). This material is distributed under the Creative Commons Attribution 4.0 International licence, http://creativecommons.org/licenses/by/4.0/, as recorded in the arXiv OAI licence metadata for this exact version and shown on the abstract page https://arxiv.org/abs/2607.15069v2. Full rights notice: licenses/arxiv-2607.15069-CC-BY-4.0.txt.\par\smallskip
\noindent\textit{Editorial scope.} Counted unit: the original \texttt{proposition} environment \texttt{prop-dim} at source lines 444--454 together with its complete proof at lines 457--479; the delivered document keeps source lines 131--480 so that every referenced label and every citation resolves inside it. Native editable source is the paper's own concatenated TeX; the AMS wrapper is editorial formatting only. No publisher class, upstream script or unrelated artwork is redistributed.
\begin{lemma}\label{lem-rankaai}
Let $A=(a_1,a_2,\ldots,a_m)\in \mathbb{R}^{n\times m}$. Let $V_A=\langle a_1,a_2,\ldots a_m \rangle$ be the linear subspace of $\mathbb{R}^{n}$ spanned by $a_i$. Let $V_A^I=\{v\in V_A| v_I=0 \}$.
Then $V_A^I \subseteq V_A$ is a linear subspace with dimension
$$\dim V_A^I=\operatorname{rank}A-\operatorname{rank}A_I.$$

In particular, if $\operatorname{rank}A=\operatorname{rank}A_I$, then $\dim V_A^I=0$.
\end{lemma}

\begin{proof}
Consider the coordinate restriction map
$$\pi_I:V_A\longrightarrow \mathbb{R}^{|I|},\qquad \pi_I(v)=v_I,$$
where $v_I$ denotes the vector consisting of the coordinates of $v$
indexed by $I$. Then $\pi_I$ is linear. Moreover,
$$\ker \pi_I=V_A^I.$$
Hence $V_A^I$ is a linear subspace of $V_A$.

Since every $v\in V_A$ can be written as $v=A c$
for some $c\in\mathbb{R}^m$, we have
$$v_I=A_I c.$$
Therefore,
$$\operatorname{Im}\pi_I=\{A_I c \mid c\in\mathbb{R}^m\},$$
which is exactly the column space of $A_I$. Hence
$$\dim \operatorname{Im}\pi_I=\operatorname{rank} A_I.$$
By the rank-nullity theorem (see, e.g., \cite[Thm 4.1.6]{artin11}),
$$\dim V_A=\dim \ker\pi_I+\dim\operatorname{Im}\pi_I.$$
Since $\dim V_A=\operatorname{rank} A$, we obtain
$$\dim V_A^I=\dim\ker\pi_I=\operatorname{rank} A-\operatorname{rank} A_I.$$
\end{proof}


\begin{lemma}\label{lem-subfull}
Suppose that $A=(a_1,a_2,\ldots,a_m)\in \mathbb{R}^{n\times m}$  and  $B=(b_1,b_2,\ldots,b_k)\in \mathbb{R}^{n\times k}$ where $a_i,~b_i\in \mathbb{R}^{n}$. Let $A_I$ and $B_I$ be the row submatrices of $A$ and $B$
indexed by $I$. Let $\langle a_1,a_2,\ldots,a_m \rangle$ and $\langle b_1,b_2,\ldots,b_k \rangle$ be linear subspaces of $\mathbb{R}^{n}$ spanned by $a_i$ and $b_i$. Let $V_A=\langle a_1,a_2,\ldots,a_m \rangle$ and 
 $V_B=\langle b_1,b_2,\ldots,b_k \rangle$. Suppose that $V_B\subseteq V_A$. Given a subset $I\subseteq [n]$, we have
$$\operatorname{rank}A_I\geq \operatorname{rank}B_I$$
and
$$\operatorname{rank}A-\operatorname{rank}A_I\geq \operatorname{rank}B-\operatorname{rank}B_I.$$
\end{lemma}
\begin{proof}
Let $a_{i,I}$ $(1\leq i\leq m)$ denote the restriction of $a_i$ to $I$ and similarly, for $b_{j,I}$ $(1\leq j\leq k)$. Then we have $A_I=(a_{1,I}, a_{2,I},\ldots,a_{m,I})$ and $B_I=(b_{1,I}, b_{2,I},\ldots,b_{k,I})$.
Since $\langle b_1,b_2,\ldots,b_k \rangle\subseteq \langle a_1,a_2,\ldots,a_m \rangle$, we have $\langle b_{1,I},b_{2,I},\ldots,
b_{k,I} \rangle\subseteq \langle a_{1,I},a_{2,I},\ldots,a_{m,I} \rangle$.
So we have
$$\operatorname{rank}B_I=\dim \langle b_{1,I},b_{2,I},\ldots,
b_{k,I} \rangle 
\leq \dim \langle a_{1,I},a_{2,I},\ldots,a_{m,I} \rangle=\operatorname{rank}A_I.$$

Let $V_A^I=\{v\in V_A| v_I=0 \}$ and $V_B^I=\{v\in V_B| v_I=0 \}$. Then we also have that $V_B^I\subseteq V_A^I$. By Lemma \ref{lem-rankaai}, we have $$ \operatorname{rank}A-\operatorname{rank}A_I=\dim V_A^I \geq \dim V_B^I=\operatorname{rank}B-\operatorname{rank}B_I.$$
\end{proof}
Then by Lemma \ref{lem-subfull} we have the following corollary.
\begin{corollary}\label{cor-subfull}
In Lemma \ref{lem-subfull},
if $\langle a_1,a_2,\ldots,a_m \rangle=\langle b_1,b_2,\ldots,b_k \rangle$ then for any subset $I\subseteq [n]$ we have
$$\operatorname{rank}A_I=\operatorname{rank}B_I$$
and  
$$\operatorname{rank}A-\operatorname{rank}A_I= \operatorname{rank}B-\operatorname{rank}B_I.$$
\end{corollary}

\subsection{Substitution by fixing a partial solution}
\


Let $F: \mathbb{R}^n\to \mathbb{R}^m$ be a polynomial mapping defined by 
$$F(x)=
\left(
\begin{array}{c}
f_1(x)\\
f_2(x)\\
\vdots\\
f_m(x)
\end{array}
\right),
$$
where $x=(x_1,x_2,\ldots,x_n)$. The corresponding system of polynomial equations $F(x)=0$ is given by
\begin{equation}\label{eq-FPE}
f_i(x_1,x_2,\ldots,x_n)=0\quad\text{where}~i=1,2,\ldots,m.
\end{equation}
The Jacobian matrix of $F$ is given by 
\begin{equation}\label{eq-jac}
J(x) =\left(\frac{\partial f_i}{\partial x_j}\right)_{i,j},~\text{where}~1\leq i\leq m,~1\leq j\leq n.
\end{equation}
Let $V(F)=\{x\in \mathbb{R}^n|F(x)=0\}$ be the solution set of (\ref{eq-FPE}). Given $s\in V(F)$, we let $\dim_s V(F)$ denote the local dimension of $s$ \cite{cox25,li25def}.  
\begin{definition}\cite{cox25}
A solution $s\in V(F)$ is called \emph{smooth} (or nonsingular) if
$\dim_s V(F)=d$, where $d=n-\operatorname{rank}J(s)$.
\end{definition}
So if $s$ is a smooth solution whose local dimension is $d$, then 
$\operatorname{rank}J(s)=n-d$. The nullspace of $J(s)$ is defined by
$$\mathrm{Null}(s)=\{x\in \mathbb{R}^n| J(s)x=0\},$$
which is a $d$-dimensional subspace of $\mathbb{R}^n$.
Choosing a (typical) basis $\{v_1,v_2,\ldots,v_d\} \subseteq~\mathrm{Null}(s)$, we define the nullspace basis matrix of $s$ by
\begin{equation}\label{eq-ns}
N(s)=(v_1,v_2,\ldots,v_d)\in \mathbb{R}^{n\times d}.
\end{equation}
Then we have $\operatorname{rank}N(s)=d$. Let $I\subseteq [n]$ be a subset and $|I|=k$. For a solution $s=(s_1,s_2,\ldots,s_n)$, we let 
$$s_I=(s_{i_1},s_{i_2},\ldots,s_{i_k}),$$
where $i_j\in I$ and $i_1<i_2<\cdots<i_k$. Then $s_I$ is called \emph{a partial solution} of $F(x)=0$ (see, e.g., \cite{cox25}). By fixing the partial solution and adjoining the constraints $x_i-s_i=0$ for $i\in I$ to $F(x)=0$, we obtain the constrained polynomial system $F_I(x)=0$:
\begin{equation}\label{eq-FI}
\left\{
\begin{array}{r@{\;}c@{\;}r@{\;}c@{\;}l}
\multicolumn{3}{r}{F(x)} & = & 0,\\
x_{i_1} & - & s_{i_1} & = & 0,\\
x_{i_2} & - & s_{i_2} & = & 0,\\
        & \vdots &     &   &   \\
x_{i_k} & - & s_{i_k} & = & 0.
\end{array}
\right.
\end{equation}
Then the corresponding polynomial mapping is:
$$F_I(x)=
\left(
\begin{array}{c}
F(x)\\
x_{i_1}-s_{i_1}\\
\vdots\\
x_{i_k}-s_{i_k}
\end{array}
\right),
$$
where $x=(x_1,x_2,\ldots,x_n)$.
Clearly, $s$ is also a solution of $F_I(x)=0$.

\begin{remark}
To solve the constrained system $F_I(x)=0$ in \eqref{eq-FI}, we
may substitute $x_i=s_i$ for all $i\in I$ into $F(x)$. This yields
a reduced polynomial system in the variables $x_j$ where $j\in[n]\setminus I$. Without of confusion, we also denote this reduced system by
$F_I(x)=0$. We refer to this procedure as \emph{making a
substitution}, or simply \emph{substitution}.
\end{remark}

Let 
\begin{equation}\label{eq-VFI}
V(F_I)=\{x\in \mathbb{R}^n|F_I(x)=0\}
\end{equation}
denote the solution set of $F_I(x)=0$. Then we have $V(F_I)\subseteq V(F)$. 

\begin{remark}
Let 
$$L_I=\{x\in \mathbb{R}^n|x_i=s_i~\text{for}~i\in I\}$$
be an affine coordinate subspace. 
So making a substitution can be understood as intersecting $V(F)$ with $L_I$. Thus, we have $V(F_I)=V(F)\cap L_I$.
\end{remark}

\begin{definition}
Let  $s\in V(F)$ be a smooth solution. A partial solution $s_I$ is called 
\emph{regular} if $s$ is a smooth solution of $F_I(x)=0$. The corresponding subset $I\subseteq [n]$ is called \emph{a regular coordinate set}.
\end{definition}
\begin{remark}
It is not hard to find a partial solution that is not regular. For example,
consider $f(x,y,z)=xy-z=0$ in $\mathbb{R}^3$, $s=(0,0,0)$ is a smooth solution. However, since (0, 0) is not a smooth point on $xy=0$, the partial solution $s_3=0$ is not regular.
\end{remark}

%Let $V(s, F)$ denote the irreducible variety of $V(F)$ which contains $s$.
%Then $$V(s,F )\cap L_I$$
%may be not irreducible. Let $V(s, F_I)\subseteq V(s,F)\cap L_I$ be the %irreducible variety that contains $s$.



Let $e_i = (0,\ldots,0,1,0,\ldots,0) \in \mathbb{R}^{1\times n}$ be the $i$-th standard basis row vector, whose $i$-th entry is $1$ and all other entries are zero. For $I=\{i_1,i_2,\ldots,i_k\}$, let
$$E_I=\left(
\begin{array}{c}
e_{i_1}\\
e_{i_2}\\
\vdots\\
e_{i_k}
\end{array}
\right).
$$
Let $J_I(x)$ be the Jacobian matrix of $F_I(x)$. It is not hard to see that 
\begin{equation}\label{eq-JI}
 J_I(x)=\left(
\begin{array}{c}
J(x)\\
E_I
\end{array}
\right).
\end{equation}

Let $\mathrm{Null}_I(x)$ denote the nullspace of $J_I(x)$. Let 
\begin{equation}\label{eq-nis}
N_I(s) 
\end{equation}
be the submatrix of $N(s)$ obtained by choosing rows indexed by $I$, and denote $d_{N_I}=\operatorname{rank}N_I(s)$.

\begin{proposition}\label{prop-nbm}
Let $s_I$ be a partial solution.
The nullspace of $J_I(s)$ is denoted by $\mathrm{Null}_I(s)$.
Then we have
$$\mathrm{Null}_I(s)=\{v|v\in \mathrm{Null}(s)~\text{and}~v_I=0\}$$
and
$$\dim \mathrm{Null}_I(s)=\operatorname{rank}N(s)-\operatorname{rank}N_I(s)=d-d_{N_I}.$$
\end{proposition}
\begin{proof}
By definition,   we have
$$\mathrm{Null}_I(s)=\{v\in \mathbb{R}^n| J_I(s)v=0\}.$$
From (\ref{eq-JI}), we can see that
if $J_I(s)v=0$, then $J(s)v=0$. So the nullspace of $J_I(s)$ is a linear subspace of the nullspace of $J(s)$.
On the other hand, if $J_I(s)v=0$, then $E_I v=0$ which implies $v_I=0$.
Conversely, if $v\in \mathrm{Null}(s)$ and $v_I=0$, then $v\in \mathrm{Null}_I(s)$.
So we have 
$$\mathrm{Null}_I(s)=\{v|v\in \mathrm{Null}(s)~\text{and}~v_I=0\}.$$

Then by Lemma \ref{lem-rankaai}, the dimension of $\mathrm{Null}_I(s)$ is
$$\dim \mathrm{Null}_I(s)=\operatorname{rank}N(s)-\operatorname{rank}N_I(s)=d-d_{N_I}.$$
\end{proof}
By Proposition \ref{prop-nbm}, we have the following corollary.
\begin{corollary}\label{cor-ddi}
Suppose that $I$ is a regular coordinate set.
If $d-d_{N_I}=0$, then $s$ is an isolated solution of $F_I(x)=0$.
If $d-d_{N_I}>0$, then $s$ is a positive-dimensional solution of $F_I(x)=0$ with local dimension $d-d_{N_I}$.
\end{corollary}


\subsection{Making the substitution by nullspace basis matrix} 
\label{subsec-first}
\

Let $F(x)=0$ be a (large) polynomial system whose solution set $V(F)$ is positive-dimensional. Assume that $s$ is a smooth point of a positive-dimensional component of $V(F)$. After making the substitution as in (\ref{eq-FI}), the original system is reduced to a smaller system $F_I(x)=0$ with solution set $V(F_I)$. If we choose $I$ such that $|I|$ is large enough, then $F_I(x)=0$ has lower complexity. We can solve it by symbolic computation software, such as, Macaulay2, Oscar.jl and Singular \cite{Dec25,Eis-M2book02,GreuelPf08}. However, after making a substitution, $s$ may become an isolated solution of $F_I(x)=0$. By Corollary \ref{cor-ddi}, we obtain an elementary first-order necessary condition for making the substitution, which can keep $s$ to be positive-dimensional in $V(F_I)$. That is, we can use the partial solution $s_I$ indexed by $I\subseteq [n]$ such that $$\operatorname{rank}N(s)-\operatorname{rank}N_I(s)>0.$$

However, when the solution set is invariant under a positive-dimensional Lie group action, such as the Brent equations\cite{Heule21,li25def}, all solutions of the reduced system may belong to a single group orbit. So in this case, the solutions we found are not new up to the group action. In the next section, we will discuss how to choose $I$ such that $V(F_I)$ contains solutions in different group orbits.

\section{Substitution modulo a Lie group action}
\label{sec-subslg}
In this section, we discuss how to choose a partial solution such
that the solution set of the reduced polynomial system contains
points from distinct group orbits. In the following discussion, the
manifold $M$ is assumed to be the smooth part of a real affine algebraic
variety in $\mathbb{R}^n$. We assume that $G$ acts linearly on the ambient space $\mathbb{R}^n$, that $M$ is $G$-invariant, and that the action of
$G$ on $M$ is obtained by restriction.
\subsection{Group orbit and tangent basis matrix} 
\

Suppose that  $M\subseteq \mathbb{R}^n$ is an $m$-dimensional manifold. Let $G$ be a $d$-dimensional Lie group that acts on $M$. Let  $\mathfrak{g}$ be the Lie algebra of $G$. Therefore, as a linear space we have $\dim \mathfrak{g}=d$. The \emph{infinitesimal generators} of the group action are given by \cite[(1.48)]{olver93li}:
\begin{equation}\label{eq-infgen}
A\cdot x =\left.\frac{d}{dt}\exp(tA)x\right|_{t=0},
\quad x \in \mathbb{R}^n,\; A \in \mathfrak{g}.
\end{equation}
So we have $A\cdot x\in \mathbb{R}^n$.


Let $\{E_1,E_2,\ldots,E_k\}\subseteq \mathfrak{g}$ be a subset.
Let $\langle E_1,E_2,\ldots,E_k\rangle$ be the linear space spanned by $\{E_1,E_2,\ldots,E_k\}$.
Let $s\in M$ be a point. Let $Gs$ be the group orbit of $s$, which is a submanifold of $M$. 
\emph{The stabilizer group} $\operatorname{Stab}_G(s)$  of $s\in M$ is defined as the subgroup that fixes $s$:
$$\operatorname{Stab}_G(s)=\{g\in G|gs=s\}.$$
The group action is \emph{semi-regular} if $\dim \operatorname{Stab}_G(s)$ is constant for all $s\in M$. The group $G$ acts \emph{regularly} if the action is semi-regular, and, in addition, for each point $s\in M$ there exist arbitrarily small neighbourhoods $U$ of $s$  with the property that each orbit of $G$ intersects $U$ in a pathwise connected subset \cite{olver93li}.

Suppose that $\langle E_1,E_2,\ldots,E_k\rangle=\mathfrak{g}$. Let $\langle E_1\cdot s, E_2\cdot s,\ldots,E_k\cdot s\rangle \subseteq \mathbb{R}^n$ denote the linear space spanned by the infinitesimal generators $\{ E_i\cdot s|i=1,2,\ldots,k\}$. Let $T_s(Gs)$ be the \emph{tangent space} of $s$ for the submanifold $Gs$. We have the following proposition.
\begin{proposition}\cite[Prop. 2.65]{olver95e}\label{prop-tdim}
 For the submanifold $Gs$, the tangent space of $s$ is given by:
$$T_s (Gs)=\mathfrak{g}\cdot s=\langle E_1\cdot s, E_2\cdot s,\ldots,E_k\cdot s\rangle.$$
For the dimension of the orbit, we have
\begin{equation*}
\dim Gs=\dim \langle E_1\cdot s, E_2\cdot s,\ldots,E_k\cdot s\rangle.
\end{equation*}
\end{proposition}

By Proposition \ref{prop-tdim}, with the notations above we have the following definition.
\begin{definition}\label{def-tsm}
The \emph{tangent basis matrix} (or infinitesimal generator coefficient matrix \cite{folver99mcf, Hubert07sa}) of $s$ with respect to the Lie group $G$ is defined as
\begin{equation}\label{eq-Ts}
T(s)=\left(E_1\cdot s, E_2\cdot s,\cdots,E_k\cdot s\right)
\in \mathbb{R}^{n\times k}.
\end{equation}
Given a subset $I\subseteq [n]$, let 
\begin{equation}\label{eq-Tis}
T_I(s)
\end{equation}
the submatrix of $T(s)$ obtained by selecting
the rows indexed by $I$. So if $|I|=r$, we have $T_I(s)\in \mathbb{R}^{r\times k}$.
\end{definition}

The following proposition gives the connection between the dimension of the stabilizer at a point and the dimension of the orbit through that point.
\begin{proposition}\cite[Prop.2.67]{olver95e}\label{prop-dimstab}
If $G$ is an $r$-dimensional Lie group acting on $M$, then the stabilizer $\operatorname{Stab}_G(x)$ of any point $x \in M$ has dimension $r-d$, where $d$ is the dimension of the orbit of $G$ through $x$.
\end{proposition}

From Proposition \ref{prop-dimstab}, we have the following corollary.
\begin{corollary}\label{cor-dimos}
Let $T(s)$ be the tangent basis matrix defined in (\ref{eq-Ts}) for a point $s\in M$. The dimension of the group orbit is given by 
$$\dim Gs= \operatorname{rank}T(s).$$
The dimension of the stabilizer is given by
$$\dim \operatorname{Stab}_G(s)=\dim \mathfrak{g}-\operatorname{rank}T(s).$$
In particular, an action is semi-regular if all the orbits have the same dimension.
\end{corollary}

\begin{definition}\label{def-tranint}
Let $K$, $L$ be two submanifolds of $M$. Suppose that the intersection $K\cap L$ is also a smooth manifold. For a point $s\in K\cap L$, we say $K\cap L$ is \emph{a local transverse intersection at} $s$, if
 $$T_s (K\cap L)=T_s K\cap T_s L.$$
If $T_s (K\cap L)=T_s K\cap T_s L$ for all $s\in K\cap L$, then we say $K\cap L$ is \emph{a transverse intersection}. 
\end{definition}
Here our definition of transverse intersection is weaker than \cite{lee12smooth}, which is called \emph{clean intersection} in \cite[Append. C.3]{Hor85III}.

Let $M\subseteq \mathbb{R}^n$ be a smooth $m$-dimensional manifold. Suppose that there is a Lie group $G$ that acts on $M$ regularly, and $\dim M>\dim G$. Then the connected components of the group orbits form a
partition of $M$. For a point $x\in M$, since $\dim M>\dim Gx$, every open neighbourhood of $x$ in $M$ intersects infinitely many group orbits.
Following the definition in \cite{lawson74fol,lee12smooth}, $M$ is endowed with \emph{a foliation}  whose \emph{leaves} consist of the connected group orbits $Gx$. That is, there exists a set of representatives $\mathcal{A}\subset M$ such that
\begin{equation*}
M=\bigsqcup_{x\in \mathcal{A}} Gx.
\end{equation*}

Let $s\in M$ be a point. Since $M\subseteq \mathbb{R}^n$, we let $s=(s_1,s_2,\ldots,s_n)$ be the coordinate of $s$. For an index set $I\subseteq [n]$, let 
\begin{equation}\label{eq-LI}
L_I=\{x \in \mathbb{R}^n\mid x_i=s_i~\text{for all}~i\in I\}  
\end{equation}
be an affine coordinate subspace. Consider the intersection 
\begin{equation}\label{eq-MI}
M_I=M\cap L_I,
\end{equation}
then we have
\begin{equation}\label{eq-sliceo}
M_I=\bigsqcup_{x\in \mathcal{A}} (Gx\cap L_I).
\end{equation}
Since $s\in Gs\cap L_I$, we have $Gs\cap L_I\neq \emptyset$ and $M_I\neq \emptyset$. In (\ref{eq-sliceo}), each connected component of $Gx\cap L_I$ is called \emph{an orbit slice}. From Definitions \ref{def-tsm} and \ref{def-tranint}, we have the following proposition.

\begin{proposition}\label{prop-dim}
Suppose that the action of $G$ on $M$ is regular.
Assume that $Gs\cap L_I$ is a connected smooth submanifold. Suppose that
$Gs$ and $L_I$ intersect transversely along $Gs\cap L_I$.
Let  $\dim_x Gs\cap L_I$ be the local dimension of $Gs\cap L_I$ at $x$. Then we have 
$$\dim_x Gs\cap L_I=\dim_s Gs\cap L_I=\operatorname{rank}T(s)-\operatorname{rank}T_I(s),\quad\text{for all}\quad x\in Gs\cap L_I.$$
The tangent space $T_s (Gs\cap L_I)$ consists of 
$$\{v|v\in \langle E_1\cdot s,\ldots, E_k \cdot s \rangle~\text{and}~ v_I=0\}.$$

In particular, if $\operatorname{rank}T_I(s)=\operatorname{rank}T(s)$, then $\dim Gs\cap L_I=0$, which implies $s$ is an isolated point of $Gs\cap L_I$.
\end{proposition}

\begin{proof}
For a smooth manifold, the local dimension is a constant which is the dimension of the manifold. That is, 
$$\dim_x Gs\cap L_I=\dim_s Gs\cap L_I,\quad\text{for all}\quad x\in Gs\cap L_I.$$
In the following, we will show
$$\dim_s Gs\cap L_I=\operatorname{rank}T(s)-\operatorname{rank}T_I(s).$$

Since $Gs\cap L_I$ is smooth and local transverse at $s$, we have 
$$\dim_s Gs\cap L_I=\dim T_s (Gs\cap L_I),$$
and
$$T_s (Gs\cap L_I)=T_s Gs \cap T_s L_I.$$

The tangent space of $L_I$ at $s$ is
$$T_s L_I=\{x\in \mathbb{R}^n|x_i=0~\text{for}~i\in I\}.$$
Suppose that $\mathfrak{g}=\langle E_1, E_2,\ldots, E_k\rangle$. Then we have
\begin{align*}
 T_s (Gs\cap L_I)=&T_s (Gs) \cap T_s L_I \\
      =&\langle E_1\cdot s,\ldots, E_k \cdot s \rangle \cap \{x\in \mathbb{R}^n|x_i=0~\text{for}~i\in I\}
\end{align*}
So $v\in T_s (Gs\cap L_I)$ if and only if $v\in \langle E_1\cdot s,\ldots, E_k \cdot s \rangle$ and $v_I=0$. Thus, by Lemma \ref{lem-rankaai} we have
$$\dim T_s (Gs\cap L_I)=\operatorname{rank}T(s)-\operatorname{rank}T_I(s).$$

So if $\operatorname{rank}T(s)-\operatorname{rank}T_I(s)=0$, then $\dim_s Gs\cap L_I=0$, which implies $s$ is isolated.
\end{proof}


\begin{thebibliography}{99}
\bibitem[Dec25]{Dec25} W.~Decker, C.~Eder, C.~Fieker, M.~Horn, and M.~Joswig, eds., \emph{The Computer Algebra System OSCAR: Algorithms and Examples}, Algorithms and Computation in Mathematics, Vol.~32, Springer, 2025.
\bibitem[Eis-M2book02]{Eis-M2book02} D.~Eisenbud, D.~R. Grayson, M.~E. Stillman, and B.~Sturmfels, eds., \emph{Computations in Algebraic Geometry with Macaulay 2}, Algorithms and Computation in Mathematics, Vol.~8, Springer-Verlag, Berlin, 2002.
\bibitem[GreuelPf08]{GreuelPf08} G.-M.~Greuel and G.~Pfister, \emph{A \textsc{Singular} Introduction to Commutative Algebra}, with contributions by O.~Bachmann, C.~Lossen, and H.~Sch\"onemann, 2nd ed., Springer, Berlin, 2008.
\bibitem[Heule21]{Heule21} Marijn J. H. Heule, M. Kauers, M. Seidl, \textit{New ways to multiply $3\times3$-matrices}, J. Symbolic Comput. \textbf{104}(2021), 899-916.
\bibitem[Hor85III]{Hor85III} L.~H\"ormander, \emph{The Analysis of Linear Partial Differential Operators III: Pseudo-Differential Operators}, Classics in Mathematics. Reprint of the 1994 edition. Springer, Berlin, 2007.
\bibitem[Hubert07sa]{Hubert07sa} E.~Hubert and I.~A. Kogan, \emph{Smooth and algebraic invariants of a group action: local and global constructions}, Found. Comput. Math. \textbf{7}(2007), 455--493.
\bibitem[artin11]{artin11} M.~Artin, \emph{Algebra}, 2nd ed., Pearson Prentice Hall, 2011.
\bibitem[cox25]{cox25} D.~A. Cox, J.~Little, and D.~O'Shea, \textit{Ideals, Varieties, and Algorithms: An Introduction to Computational Algebraic Geometry and Commutative Algebra}, 5th ed., Undergraduate Texts in Mathematics, Springer, Cham, 2025.
\bibitem[folver99mcf]{folver99mcf} M.~Fels and P.~J. Olver, \emph{Moving Coframes: II. Regularization and Theoretical Foundations}, Acta Applicandae Mathematicae \textbf{55}(1999), 127--208. %
\bibitem[lawson74fol]{lawson74fol} H.~B.~Lawson, Jr., \emph{Foliations}, Bull. Amer. Math. Soc. \textbf{80}(1974), no.~3, 369--418.
\bibitem[lee12smooth]{lee12smooth} J.~M.~Lee, \emph{Introduction to Smooth Manifolds}, 2nd ed., Grad. Texts in Math., Vol.~218, Springer, 2012.
\bibitem[li25def]{li25def} X.~Li, L.~Zhang, and Y.~Ke, \emph{Deflation Conjecture and Local Dimensions of Brent Equations}, Exp. Math. \textbf{34}(2025), no.~3, 488--501.
\bibitem[olver93li]{olver93li} P.~J. Olver, \textit{Applications of Lie Groups to Differential Equations}, 2nd ed., Grad. Texts in Math., Vol.~107, Springer, New York, 1993.
\bibitem[olver95e]{olver95e} P.~J. Olver, \textit{Equivalence, Invariants and Symmetry}, Cambridge Univ. Press, Cambridge, 1995.
\end{thebibliography}
\end{document}
